\chapter*{Glossary}
\addcontentsline{toc}{chapter}{Glossary}

\begin{description}

  \item[Monosemanticity] A property of a neuron or feature whose activation corresponds to a single, identifiable concept or pattern, in contrast to polysemantic units that respond to multiple unrelated concepts due to superposition.

  \item[Circuit] A subgraph of the model's computational graph comprising a small set of
    components (attention heads, MLP layers, or individual features) that together implement a
    specific behaviour.

  \item[Grokking] The phenomenon whereby a neural network trained on a small dataset first
    memorises the training examples and
    then, after many additional gradient steps, abruptly transitions to a generalising solution.

  \item[SAE --- Sparse Autoencoder] A dictionary-learning module trained to decompose
    neural network activations into a sparse linear combination of learned feature directions.

  \item[Unembedding matrix] The learned linear map that projects 
  the final residual-stream vector onto vocabulary logits, 
  from which the next-token distribution is obtained via softmax.

  \item[Representation engineering] A methodology for extracting and analysing concept
    directions in the activation space of a neural network by computing the mean difference
    in hidden-state vectors between semantically contrasting inputs.

  \item[Activation patching] A causal intervention technique in which a single intermediate
    representation, here the residual-stream vector at a chosen layer and token position,
    is replaced by the corresponding value from a different input.
    The resulting change in output logits quantifies the causal sufficiency of that
    representation for driving the model's prediction.

  \item[Crossover] In an intervention experiment, the point at which the source-condition token
    probability falls below the target-condition token probability as intervention strength
    increases. 
    
  \item[Anchor position] A designated token position $t^*$ in the prompt at
    which residual-stream vectors are collected.  The choice of anchor determines which part
    of the prompt's processing is examined, as different anchors can reveal
    different aspects of model internals.

\end{description}
